% ============================================================================
% 08_discussion.tex — Discussion (sec:discussion)
% Per plan/plan.md §3 row 8: design implications + practitioner box (Design-C
% graft) + scope conditions + one reconciled cost table (tab:cost, mined from
% writeup/contents/limitation.tex — that file IS the old Discussion) +
% limitations. Merge pass: scope conditions rewritten to the genuine
% gpt-5-mini cells (constraints migrate, not vanish); traces double-
% dissociation paragraph added; practitioner box and limitations updated to
% the invariance-content story and the rescoped DA observability accounting.
% ============================================================================

\section{Discussion}\label{sec:discussion}

\subsection{Design Implications}\label{sec:discussion-implications}

The ranking's most actionable feature is that a nearly costless lever sits close to its top. A one- or two-sentence honest description of a mechanism's safety invariant recovers a large share of what a full obviously strategy-proof redesign buys, without touching the allocation rule, the message space, or the extensive form: Payoff Safety cut the cross-model auction deviation from $-\$2.67$ to $-\$1.53$, and the clock-framing description substantially reduced SMAD in three of the four model families while backfiring in the fourth (pooled $\rho=+0.30$, wide interval; Section~\ref{sec:ranking-descriptions}). The supplementary matching experiment pushes the same point further: Rejection Safety cut direct-DA error from $4.2\%$ to $0.2\%$, and a plain textbook statement of strategy-proofness drove it to $0.0\%$ in all four models (Appendix~\ref{app:da}). For a platform whose participants increasingly include LLM agents, these are deployable as interface text or API documentation at approximately zero engineering cost. They are also honest in the sense of \citet{gonczarowski2023strategyproofness}: they state true properties of the mechanism rather than steering behavior through misdirection, so a designer can adopt them without incurring the credibility costs of manipulative framing. The comparison with restatements disciplines the recommendation, and more sharply than a simple placebo: it is not ``more explanation'' that helps. Menu restatements of the auction are null on average and sign-mixed across models (Section~\ref{sec:ranking-descriptions}), and in the matching domain the same menu format splits by content---a menu that includes the invariance statement improves play like a safety description, while a menu that merely restates the outcome computation makes play significantly \emph{worse} (Appendix~\ref{app:da}). The active ingredient is the invariance content, not the explanation.

As LLMs increasingly participate in markets, the design of the mechanisms they interact with will matter. Our results suggest that simple mechanisms, and clear (and honest) descriptions of mechanism properties, can serve as a form of alignment: channeling artificial agents toward the outcomes the mechanism was designed to produce, without requiring that they fully understand why. And because every lever with a human anchor moves human participants in the same direction it moves LLM agents (Section~\ref{sec:humans}), a designer serving a mixed human--AI market need not choose between populations: simplicity, on this evidence, is a common good across participant types.

The reasoning traces add a caution about \emph{how} these text levers work---and about how their efficacy should be audited. The dissociation of Section~\ref{sec:ranking-traces} shows that the effective levers change behavior without changing stated reasoning: Payoff Safety cuts the mean absolute deviation from $\$3.20$ to $\$1.95$ while its invariant is echoed in none of the six hundred treated traces and the rate of correct payment-rule statements is flat, so a classical mediation analysis fails at the first stage and bounds any comprehension-mediated share of the effect at roughly one third; the payoff tree behaves the same way (bids move, rule-rehearsal language does not). Conversely, the worst-case scaffold visibly changes what models say---worst-case language jumps from $2\%$ to $43\%$ of traces---without moving bids, and the menu restatement moves neither: across every lever tested, language and bids never move together. The channel through which descriptions operate is therefore behavioral rather than deliberative: the text changes what agents \emph{do} without passing through what they \emph{articulate}. For practitioners this cuts in both directions. A designer cannot certify a description by checking that agents verbalize the right principle (the best lever produces no such verbalization), and cannot dismiss one because comprehension checks fail; symmetrically, a passed manipulation check on stated reasoning is no evidence that behavior has moved. Interventions on LLM participants should be evaluated on realized play, not on elicited understanding.

The backfiring tier carries an equally practical warning. Forward-planning and belief-formation scaffolds are not exotic: they are precisely the kind of ``think ahead, model the other side'' instructions that agent frameworks and prompt libraries routinely inject into deployed LLM pipelines \citep{fish2024algorithmic}. In strategy-proof environments, our evidence says such instructions are at best inert and at worst substantial liabilities, with the damage concentrated in the weakest models (Section~\ref{sec:ranking-scaffolds}). We summarize the operational content of the ranking as follows.

\begin{center}
\fbox{\begin{minipage}{0.93\linewidth}
\vspace{0.4em}
\textbf{Practitioner guidance: fielding strategy-proof mechanisms to LLM participants.}
\begin{itemize}[leftmargin=1.4em, itemsep=2pt, topsep=3pt]
    \item \emph{Prefer OSP implementations where available.} The ascending-clock format produced the best play in every model we measured, consistent with its record in human markets \citep{li2017obviously, milgrom2020clock}.
    \item \emph{If you must keep a direct mechanism, state its safety invariant---do not merely restate its mechanics.} Add a safety-exposing sentence (``your bid determines \emph{whether} you win, not \emph{what} you pay''; ``rejections only redirect, never eliminate'') and, where the interface permits, a payoff-tree explanation of what happens under each outcome. Each is a one-paragraph text change; each independently closed roughly half the gap to the OSP benchmark in auctions, and the safety description nearly closed it in DA. Mechanical restatements of the outcome mapping are useless-to-harmful: menu-style restatements were null on average and sign-mixed in auctions, and a menu that restated the computation \emph{without} the invariance statement significantly worsened matching play. The invariance content is the active ingredient.
    \item \emph{Never prompt agents to plan ahead or to model opponents.} Lookahead and belief-formation scaffolds worsened play in nearly every configuration tested---most severely for the weakest models---by activating strategizing that the mechanism makes unnecessary.
    \item \emph{Audit vendor ``reasoning enhancements'' as potential backfire levers.} Agent frameworks that silently inject planning steps or opponent modeling into prompts are, in our taxonomy, Family-C interventions of the backfiring kind; benchmark them against truthful play before deployment.
\end{itemize}
\vspace{0.2em}
\end{minipage}}
\end{center}

\subsection{Scope Conditions}\label{sec:discussion-scope}

The ranking is a statement about a bounded-rationality regime, and both of its boundaries are measured rather than assumed. The upper boundary is now measured with the full lever grid rather than inferred from baseline play alone: we ran the complete intervention battery on four frontier reasoning models (gpt-5-mini, gpt-5, Claude Sonnet 5, Gemini 2.5 Flash; Section~\ref{sec:ranking-synthesis}, Table~\ref{tab:frontier-grid}), and the picture is three-sided. First, the ranking's positive rungs \emph{compress to nothing}: all four models bid essentially exactly truthfully at every sealed-bid baseline (mean deviations $0.00$--$0.30$, SMAD $0$--$2.3\%$), so safety descriptions, the payoff tree, and belief scaffolds leave play unchanged---the preservation index $\rho$ of Section~\ref{sec:metrics}, a ratio to the baseline deviation, becomes uninformative as $D(\text{baseline}) \to 0$. Second, the constraints \emph{migrate} rather than vanish across formats: gpt-5-mini plays second price and (correctly shaded) first price essentially perfectly, yet still fails the harder third-price format outright, bidding roughly its value where the equilibrium calls for bidding above it (three-bidder mean deviation $+\$14.76$; SMAD $60\%$; Appendix~\ref{app:ablations}), so wherever equilibrium reasoning is genuinely required a bounded-rationality regime persists even at the frontier. Third, and the reason compression is not the whole story, \emph{the backfire rungs keep their teeth}: the menu restatement collapses Claude Sonnet 5 from exact truthfulness to a $35.5\%$ deviation ($d=-3.8$, $p<0.001$) and degrades Gemini 2.5 Flash, and one sentence of affiliation language in a clock description triggers a misapplied winner's-curse script in the same model (SMAD $0.3\% \to 8.6\%$), removed entirely by an IPV-worded rerun (Section~\ref{sec:ranking-synthesis}). The frontier failure mode is thus \emph{retrieval selection}: surface wording chooses which memorized playbook runs, and a wrong match is executed flawlessly. A presentation choice with no incentive content can still break an otherwise perfect strategic reasoner, so the practitioner guidance above remains binding even where the simplifying levers have nothing left to preserve. Our main model grid deliberately comprises well-established, widely deployed model families rather than frontier reasoning models: the design question requires agents capable enough to parse the rules of a mechanism yet still susceptible to the reasoning failures that simple design aims to remedy. Deployed fleets will occupy this regime for the foreseeable future: high-volume agentic transactions run on small, cheap models for cost and latency reasons, so the population of economic agents that mechanisms will actually face looks more like our grid than like the frontier. At the lower boundary, Llama-3-8B posts a mean SMAD of $57.58\%$---worse than the reconstructed human benchmark in six of seven formats (Appendix~\ref{app:ablations})---indicating a capability floor below which agents cannot reliably parse the mechanism at all, and beneath which no descriptive lever can be expected to help. The ranking applies to the wide band between these boundaries.

\subsection{What the Testbed Costs}\label{sec:discussion-cost}

The traditional instruments for questions like ours are incentivized human experiments and observational field data, and both are expensive ways to search a large design space: the laboratory test of obvious strategy-proofness engaged 548 subjects across a main experiment and a one-shot follow-up, with documented payments on the order of \$5{,}400 for the main experiment alone \citep{li2017obviously}; incentivized description experiments recruited 542 subjects with average payments of \$15--26 per subject \citep{gonczarowski2024describing}; and observational platform data require controls for seller and bidder heterogeneity that complicate the causal inference a lever ranking needs \citep{manning2024automated}. The synthetic testbed sits at the exploratory end of this pipeline, extending the cost reduction that online laboratories brought to behavioral economics \citep{horton2011online}: the full battery reported in this paper---more than 1{,}000 auctions and 5{,}000 individual agent decisions across two domains, four model families, and the complete lever grid---cost under \$400 in API fees, and a single treatment cell of 50 repetitions costs roughly \$10 (Table~\ref{tab:cost}). These economics are most useful read jointly with Section~\ref{sec:humans}: cheap synthetic screening ranks the candidate levers; expensive human experiments are reserved for the rungs where the ranking makes novel, falsifiable predictions.

\begin{table}[t]
\centering
\small
\resizebox{\textwidth}{!}{%
\begin{tabular}{@{}llll@{}}
\toprule
\textbf{Instrument} & \textbf{Population} & \textbf{Scale} & \textbf{Approx.\ cost} \\
\midrule
OSP laboratory experiment \citep{li2017obviously} & Human & 548 subjects (144 main $+$ 404 follow-up) & $\approx\$5{,}400^{\dagger}$ \\
Description experiments \citep{gonczarowski2024describing} & Human & 542 subjects, incentivized lab & \$15--26/subject$^{\ddagger}$ \\
This paper: full lever battery & LLM ($4$ families) & $>$1{,}000 auctions; $>$5{,}000 decisions & $<\$400$ \\
This paper: one treatment cell & LLM & $50$ repetitions & $\sim\$10$ \\
\bottomrule
\end{tabular}}
\caption{Approximate cost of instruments for testing simplicity-preserving design levers. $^{\dagger}$\citet{li2017obviously}'s main experiment paid 144 subjects an average of \$37.54 (including a \$20 show-up fee), $\approx\$5{,}400$ total; documented for the main experiment only. The paper's one-shot follow-up added 404 subjects (\$5 show-up plus incentive payments), for 548 subjects overall, but does not report an average payment for the follow-up, so a total cost across both experiments is not recoverable from the published figures. $^{\ddagger}$\citet{gonczarowski2024describing} paid Prolific participants an average of £15.0 ($\approx\$19$) and Cornell participants an average of \$25.5 across $542$ subjects (Appendix A.1 of that paper).}
\label{tab:cost}
\end{table}

\subsection{Limitations}\label{sec:discussion-limitations}

\paragraph{The human reconstruction is model-based.} Our human benchmarks are moment-matched mixture reconstructions of published summary statistics, not raw experimental microdata (Section~\ref{sec:reconstruction}, Appendix~\ref{app:human}). We therefore confine human--LLM comparisons to rankings, directions, and comparative statics, attach no significance tests to reconstructed levels, and note that the ranking itself (Section~\ref{sec:ranking}) never conditions on the reconstructions---they anchor transferability only. \TODO{Reconstruction-uncertainty bootstrap bands on every human anchor (Appendix~\ref{app:human}) pending.}

\paragraph{Prompt-template specificity.} Each lever is instantiated as one template per domain, so the estimated effects are, strictly, effects of particular wordings. Three features argue against a prompt-engineering artifact: the taxonomy was fixed ex ante from published theory (Section~\ref{sec:framework}), the menu restatement functions as a built-in restatement control (null on average, though sign-mixed across models; Section~\ref{sec:ranking-descriptions}), and half the taxonomy backfires---a pattern opportunistic prompt search would not produce and would not publish. But invariance to paraphrase remains an assumption rather than a measurement. \TODO{Paraphrase robustness: three paraphrases $\times$ two headline levers (plan \S6).}

\paragraph{Direction shares are convention-dependent.} Statements about the \emph{direction} of deviations (overbidding versus underbidding shares) depend on the classification convention---whether shares are taken over all bids or only over bids that deviate from value, and how a tolerance band around the benchmark is drawn. We use the repo-canonical convention throughout Section~\ref{sec:humans} (share of bids strictly above value among all bids; the human SPSB overbid share is $67.2\%$, or $92.2\%$ among deviating bids), while the directional decomposition of Section~\ref{sec:validation} removes a $\pm2\%$ band before classifying. The qualitative reversal---humans overbid, LLMs underbid---is robust to the choice, but individual shares should not be compared across conventions.

\paragraph{Model-version dependence.} All results are estimated on pinned model snapshots (Section~\ref{sec:design}); provider updates, changes to default system prompts, and fleet turnover can shift behavior in ways we cannot forecast. The ranking should be read as a measurement protocol to be re-estimated as agent populations evolve, not as a constant of nature---and the frontier migration documented above, with dominant-strategy formats solved first and equilibrium-reasoning formats persisting, already indicates the likely direction of drift.
